\section{Task-Pipeline-Analyse}
\label{sec:results_pipeline}

In diesem Abschnitt wird der Fortschritt der Probanden entlang der Aufgaben-Pipeline (T1--T5) dargestellt. Berichtete werden Häufigkeiten für Intention (Chat), Umsetzung (Code-Artefakte) und den Statusausweis \textit{Execution observed} (Integrationsläufe).

\subsection{Methodische Definitionen}
Für die Auswertung werden drei Statuskategorien verwendet:
\begin{itemize}
    \item  \textbf{Attempted (Versucht):} Ein Task gilt als  \textit{Attempted}, wenn er im Chat-Verlauf explizit als Ziel, Anforderung oder nächster Arbeitsschritt formuliert wurde. Die bloße Existenz von Dateien ohne entsprechenden Kontext im Chat genügt nicht.
    \item  \textbf{Implemented\textsubscript{static} (Umgesetzt, Artefakt/Heuristik):}\\
    Ein Task gilt als  \textit{Implemented\textsubscript{static}}, wenn er in den Auswertungsartefakten als implementiert markiert ist (heuristische, statische Indikatoren; z.\,B.\ Feld \texttt{implemented}\\
    in \path{evaluation/task_pipeline_v3.json}\\
    bzw.\ Implementationsklassen in \mbox{\path{evaluation/deep_code_analysis.json}}).
    \item  \textbf{Execution observed (Statusausweis):} Ein Task gilt als  \textit{Execution observed}, wenn die zugehörigen automatisierten Integrationsläufe erfolgreich durchlaufen wurden (Feld \texttt{verified} in \path{evaluation/task_pipeline_v3.json}).
\end{itemize}

\noindent \textbf{Abgeleitete Evidenzstufe.} In Kapitel~6 wird eine abgeleitete Evidenzstufe\\
\BeginAccSupp{unicode,method=pdfstringdef,ActualText={extitImplemented* verwendet:}}
\csname textit\endcsname{Implemented$^\ast$} verwendet:
\EndAccSupp{}
\[
\begin{aligned}
    \csname textit\endcsname{Implemented$^\ast$} &:=  \textit{Attempted} \wedge (\textit{Implemented\textsubscript{static}} \vee  \textit{Execution observed}).
\end{aligned}
\]
Damit gilt per Konstruktion:  \textit{Execution observed} $\subseteq$  \textit{Implemented$^\ast$}.  \textit{Execution observed} ist jedoch nicht notwendigerweise ein Subset von  \textit{Implemented\textsubscript{static}}:  \textit{Implemented\textsubscript{static}} basiert auf heuristischen, strukturellen Erkennungsregeln; unter den gegebenen Rahmenbedingungen ist es möglich, dass funktionsfähige Implementationen nicht als  \textit{Implemented\textsubscript{static}} erkannt werden (z.\,B. abweichende Benennungen, alternative Architekturpfade oder Implementationen außerhalb der erwarteten Muster).  \textit{Execution observed} ist ein methodengebundenes Diagnose-Signal;  \textit{Implemented\textsubscript{static}} ist ein Artefakt-Signal.

\noindent \textbf{Hinweis zur Ausführungsmethodik.} Für die Pipeline-Auswertung wurden die automatisierten Integrationsläufe (insbesondere T2, sowie in Teilen T3--T5) als API-Integrationsläufe (\textit{Supertest} gegen Express) gestaltet. \textit{Execution observed} bedeutet in dieser Studie explizit: „durchläuft die aktuellen Integrationsläufe“ \textendash{} nicht „fachlich vollständig korrekt gelöst“.

\subsection{Pipeline-Ergebnisse (T1--T5)}
\begin{table}[ht]
\centering
\begingroup
\newcommand{\NoExtract}[1]{\BeginAccSupp{method=escape,ActualText={}}#1\EndAccSupp{}}
\newcommand{\ExtractOnly}[1]{\smash{\rlap{\BeginAccSupp{unicode,method=pdfstringdef,ActualText={#1}}\textcolor{white}{.}\EndAccSupp{}}}}
\def\TaskPipelineTableActualText{Task^^JAttempted (𝑛) Implemented* (𝑛) Execution observed (𝑛)^^JT1: Persistenz^^J18 (100\%)^^J17 (94\%)^^J10 (56\%)^^JT2: Auth^^J18 (100\%)^^J17 (94\%)^^J4 (22\%)^^JT3: Kategorien^^J15 (83\%)^^J13 (72\%)^^J13 (72\%)^^JT4: Attachments^^J13 (72\%)^^J8 (44\%)^^J0 (0\%)^^JT5: Kalender^^J13 (72\%)^^J0 (0\%)^^J0 (0\%)^^J^^JTabelle 6.2: Task-Pipeline: Von der Intention zur testgebundenen Bestätigung (𝑁 = 18,^^J𝑛 = Anzahl Branches).^^JextitAttempted stammt aus Chat-Signalen. Implemented* ist definiert als: Attempted und^^J(Implementedstatic oder Execution observed). Execution observed stammt aus V3-Integrationsläufen^^J(V3 bezeichnet die für diese Studie definierte Integrationslauf-Konfiguration).}
\NoExtract{\begin{tabular}{|l|c|c|c|}
\hline
    \csname textbf\endcsname{Task} &  \textbf{Attempted ($n$)} &  \textbf{Implemented$^\ast$ ($n$)} &  \textbf{Execution observed ($n$)} \\
\hline
T1: Persistenz & 18 (100\%) & 17 (94\%) & 10 (56\%) \\
T2: Auth & 18 (100\%) & 17 (94\%) & 4 (22\%) \\
T3: Kategorien & 15 (83\%) & 13 (72\%) & 13 (72\%) \\
T4: Attachments & 13 (72\%) & 8 (44\%) & 0 (0\%) \\
T5: Kalender & 13 (72\%) & 0 (0\%) & 0 (0\%) \\
\hline
\end{tabular}}
\ExtractOnly{\TaskPipelineTableActualText}
\NoExtract{\caption{Task-Pipeline: Von der Intention zur testgebundenen Bestätigung ($N = 18$, $n$ = Anzahl Branches). \textit{Attempted} stammt aus Chat-Signalen. \textit{Implemented$^\ast$} ist definiert als: \textit{Attempted} und (\textit{Implemented\textsubscript{static}} oder \textit{Execution observed}). \textit{Execution observed} stammt aus V3-Integrationsläufen (V3 bezeichnet die für diese Studie definierte Integrationslauf-Konfiguration).}}
\label{tab:task_pipeline_v2}
\endgroup
\end{table}

\noindent \textbf{Messung.} Berichtete werden je Task (T1--T5) die Häufigkeiten der Pipeline-Status \textit{Attempted} (Chat-Signal), \textit{Implemented$^\ast$} (kombinierte Evidenzstufe) und \textit{Execution observed} (methodengebundenes Signal) über $N=18$ Branches.
\noindent \textbf{Hauptbefund.} Für T1/T2 liegt \textit{Attempted} jeweils in 18/18 Branches vor, während \textit{Execution observed} für T4/T5 in 0/18 Branches erreicht wird.
\noindent \textbf{RQ-Bezug.} Die Statushäufigkeiten bilden die deskriptive Ergebnisbasis für RQ1 und werden in Kapitel~7 für die Einordnung der Befunde und die Pipeline-Diskussion verwendet.

\begin{figure}[ht]
\centering
\makebox[\linewidth][c]{%
\begin{tikzpicture}
\begin{axis}[
    xbar,
    width=0.92\linewidth,
    height=7.2cm,
    xmin=0,
    xmax=18,
    xtick={0,2,...,18},
    scaled x ticks=false,
    tick label style={/pgf/number format/fixed,/pgf/number format/precision=0},
    xticklabel style={font=\scriptsize},
    xlabel={Anzahl Branches ($N=18$)},
    xlabel style={font=\small},
    symbolic y coords={T1: Persistenz,T2: Auth,T3: Kategorien,T4: Attachments,T5: Kalender},
    ytick=data,
    y dir=reverse,
    bar width=6pt,
    enlarge y limits=0.45,
    trim axis left,
    trim axis right,
    axis lines=left,
    xmajorgrids,
    grid style={black!8, line width=0.15pt},
    tick style={black!50},
    legend style={
        at={(0.5,-0.18)},
        anchor=north,
        legend columns=3,
        font=\small,
        draw=none
    },
    point meta=explicit symbolic,
    nodes near coords,
    nodes near coords style={
        font=\scriptsize,
        anchor=west,
        xshift=3pt,
        text=black!80
    },
    clip=false,
]
    % Attempted (hellgrau)
    \addplot[fill=black!15, draw=black!55, line width=0.2pt, bar shift=9pt]
    coordinates {(18,T1: Persistenz) [18] (18,T2: Auth) [18] (15,T3: Kategorien) [15] (13,T4: Attachments) [13] (13,T5: Kalender) [13]};

    % Implemented* (mittelgrau, visuell dominanter)
    \addplot[fill=black!55, draw=black, line width=0.35pt, bar width=7pt, bar shift=0pt]
    coordinates {(17,T1: Persistenz) [17] (17,T2: Auth) [17] (13,T3: Kategorien) [13] (8,T4: Attachments) [8] (0,T5: Kalender) [{}]};

    % Execution observed (dunkelgrau/schwarz; zusätzlich Schraffur)
    \addplot[fill=black!85, draw=black, line width=0.4pt, pattern=north east lines, pattern color=black, bar shift=-9pt]
    coordinates {(10,T1: Persistenz) [10] (4,T2: Auth) [4] (13,T3: Kategorien) [13] (0,T4: Attachments) [{}] (0,T5: Kalender) [{}]};

    % 0-Werte (ohne Textlabel) optional als Marker am Achsenstart
    \addplot[only marks, mark=o, mark size=2.6pt, draw=black, mark layer=foreground, forget plot, bar shift=0pt]
    coordinates {(0,T5: Kalender)};
    \addplot[only marks, mark=x, mark size=3.0pt, draw=black, mark layer=foreground, forget plot, bar shift=-9pt]
    coordinates {(0,T4: Attachments) (0,T5: Kalender)};

    \legend{Attempted, Implemented$^\ast$, Execution observed}
\end{axis}
\end{tikzpicture}
}%
\caption{Pipeline-Ergebnisse für T1--T5 ($N=18$) als gruppierte Balken:  \textit{Attempted} (Chat-Signal),  \textit{Implemented$^\ast$} (definiert als:  \textit{Attempted} und (\textit{Implemented\textsubscript{static}} oder \textit{Execution observed})) und  \textit{Execution observed} (testgebunden: bestandene V3-Integrationsläufe). V3 bezeichnet die für diese Studie definierte Integrationslauf-Konfiguration; die Werte entsprechen Tabelle~\ref{tab:task_pipeline_v2}.}
\label{fig:task_pipeline_funnel_t1_t5}
\end{figure}

\noindent \textbf{Messung.} Visualisiert werden die in Tabelle~\ref{tab:task_pipeline_v2} berichteten Statushäufigkeiten (\textit{Attempted}, \textit{Implemented$^\ast$}, \textit{Execution observed}) als Balken pro Task.
\noindent \textbf{Hauptbefund.} Über T1--T5 liegen \textit{Implemented$^\ast$} und \textit{Execution observed} unter \textit{Attempted} und erreichen bei T4/T5 den Wert 0.
\noindent \textbf{RQ-Bezug.} Die Abbildung ist eine komprimierte Darstellung der Pipeline-Ergebnisse und ergänzt die Ergebniszusammenfassung zu RQ1 in Kapitel~7.

\subsection{Kurzbefunde pro Task (T1--T5)}
\noindent \textbf{T1 (Persistenz).} T1 ist in 18/18 Branches  \textit{Attempted}, in 17/18 Branches  \textit{Implemented$^\ast$} und in 10/18 Branches  \textit{Execution observed}.

\noindent \textbf{T2 (User Accounts/Auth).} T2 ist in 18/18 Branches  \textit{Attempted}, in 17/18 Branches  \textit{Implemented$^\ast$} und in 4/18 Branches  \textit{Execution observed}.

\noindent \textbf{T3 (Kategorien).} T3 ist in 15/18 Branches  \textit{Attempted}, in 13/18 Branches  \textit{Implemented$^\ast$} und in 13/18 Branches  \textit{Execution observed}. In \path{evaluation/task_pipeline_v3.json} ist T3 dabei in 7/18 Branches als \texttt{implemented} (\textit{Implemented\textsubscript{static}}) markiert.

\noindent \textbf{T4 (Attachments).} T4 ist in 13/18 Branches  \textit{Attempted}, in 8/18 Branches  \textit{Implemented$^\ast$} und in 0/18 Branches  \textit{Execution observed}.

\noindent \textbf{T5 (Kalender/Multi-Day).} T5 ist in 13/18 Branches  \textit{Attempted}, in 0/18 Branches  \textit{Implemented$^\ast$} und in 0/18 Branches  \textit{Execution observed}.

% DATENQUELLEN:
% Chat-Logs (für "Attempted"-Status)
% Code-Analyse (für "Implemented"-Status)
% Test-Reports (für "Execution observed"-Status)
%
% HINWEISE:
% Definition "Attempted": Explizite Erwähnung im Chat
% Definition "Implemented": Code vorhanden
% Definition "Execution observed": Tests bestanden
% Fokus auf die Statushäufigkeiten über die Tasks hinweg
